\subsection{Построение конвейера системы идентификации объектов в реальном времени}
По умолчанию, возможность разбора выходных тензоров Yolo в Nvidia Deepstream не предусмотрена. Однако, мы нашли решение, которое позволило использовать Nvidia Deepstream совместно с Nvidia Triton. При этом, Nvidia Deepstream занимается работой с видеопотоком, а анализ кадров передается в Nvidia Triton, где выполняется работа нейронных сетей. Полученные результаты в виде выходных тензоров модели возвращаются в Deepstream, где происходит их разбор и привязка к соответствующим кадрам в потоке.

Есть несколько способов, которые можно использовать для решения проблемы разбора выходных тензоров Yolo. 

Первый из таких способов - это реализация использования Yolo в Triton в качестве Python кода. При этом на вход модели будут поступать кадры, которые затем будут исполняться в модели и проходить через подавление немаксимумов для выходного тензора. После этого результат будет возвращен в виде вектора, который может быть разобран в формате, который Deepstream сможет распознать, например Faster R-CNN. Однако, следует отметить, что этот подход имеет низкую производительность по сравнению с другими способами. Это происходит потому, что выполнение модели и подавление немаксимумов в python\_backend Nvidia Triton медленнее, чем в других вариантах. 

Второй это рализовать свой разбор выходного тензора Yolo на C++ и подключить его в Nvidia Deepstream как динамическую библиотеку. Для этого нужно будет описать две функции класса для преобразования вектора перед передачей в Nvidia Triton и преобразование выходного тензора, где нужно описать ограничивающий прямоугольник для каждого объекта и прикрепить его к кадру. Этот способ быстрее первого, так как мы можем использовать модель непосредственно соотвествующим сервером, например pytorch\_server для Pytorch модели. Также мы можем конвертировать и использовать модель в другом формате, ONNX или TensorRT, что даст еще большую производительность. Минусом данного способа является необходимость реализации подавления немаксимумов на С++ и применять его при преобразовании выходного тензора, это означает, что мы не сможем обеспечить совместимость между YoloV3-YoloV7 и YoloV8, так как они имеют разный формат выходных тензоров и  подавления немаксимумов придется делать используя процессор. Можно остановиться на этом способе, но я нашел, как можно решить проблемы, которые я описал выше.

Третий способ основан на втором, в нем мы будем реализовывать свой разбор выходного тензора Yolo на C++ и подключать его в Nvidia Deepstream как динамическую библиотеку. Различие лишь в том, что мы на этапе конвертации модели добавим подавления немаксимумов в саму структуру сети. Это позволит унифицировать формат вывода YoloV8 и предыдущих версий, а также производить подавления немаксимумов на видеокарте, что ускорит выполнение.

Я буду использовать третий способ, рассмотрим его подробнее на примере YoloV8. Начнем с конвертации и добавления подавления немаксимумов. Для этого сначала нужно конвертировать модель из формата Pytorch в формат ONNX. ONNX (Open Neural Network Exchange) - это открытый формат обмена моделями глубокого обучения между различными фреймворками и инструментами машинного обучения. Это будет промежуточной ступенью для конвертации модели TensorRT. Напишем простой код для конвертации Yolo в ONNX. Пример приведен на рисунке ~\ref{src:src2}.

\begin{figure}
\begin{lstlisting}[language=Python]
from ultralytics import YOLO
model = YOLO("yolov8n.pt")
success = model.export(format="onnx", batch=2)
\end{lstlisting}
\caption{Конвертация Yolo в ONNX}
\label{src:src2}
\end{figure}

Далее нам нужно окружение с установленным TensorRT, я использовал для этого docker контейнер deepstream:6.2-triton, так как там уже установлены эти компоненты. Для добавления подавления немаксимумов и конвертации в TensorRT я использовал git репозиторий YOLO Series TensorRT \cite{yolotrt2022}, для конвертации этот репозиторий использует trt-exec, который предустановлен в нашем контейнере. Также для его работы нам необходимо необходимо установить еще несколько python зависимостей. Скрипт для установки зависимостей приведен на рисунке ~\ref{src:src2_1}.

\begin{figure}
\begin{lstlisting}
pip install --upgrade setuptools pip --user
pip install nvidia-pyindex
pip install --upgrade nvidia-tensorrt
pip install pycuda
pip install Pillow
\end{lstlisting}
\caption{Установка зависимостей}
\label{src:src2_1}
\end{figure}

Теперь установив все зависимости и конвертировав модель в ONNX мы можем приступить к конвертации в TensorRT. Для этого нужно вызвать скрипт для конвертации. 
python export.py -o yolov8n.onnx -e yolov8n.trt --end2end --v8

Разберем параметры параметры, представленные на рисунке ~\ref{src:src2_2}.
\begin{figure}
\begin{lstlisting}
-o — путь до ONNX файла
-e — путь до выходного TensorRT
--v8 — параметр, который нужно указать при работе с Yolov8
--precision [fp32 / fp16 / int8] — точность весов модели (по умолчанию fp16)
--end2end — параметр отвечающий за добавления подавления немаксимумов
--max_det — максимальное количество объектов в кадре после подавления немаксимумов (по умолчанию 100)
--conf_thres —порог достоверности (по умолчанию 0.4)
--iou_thres — порог пересечений и объединений (по умолчанию 0.5)
\end{lstlisting}
\caption{Параметры конвертера}
\label{src:src2_2}
\end{figure}

На выходе мы получаем получаем модель в формате TensorRT с несколькими выходными тензорами, параметры которой представлены на рисунке ~\ref{src:src2_3}.

\begin{figure}
\begin{lstlisting}
boxes: (batch_size, max_det, 4) — где координаты формата XYXY
classes: (batch_size, max_det) — классы объектов
num: (batch_size,1) — количество объектов на кадре
scores: (batch_size, max_det) — достоверность для каждого объекта
\end{lstlisting}
\caption{Выходные тензоры}
\label{src:src2_3}
\end{figure}

Зная это описываем файл config.pbtxt, где указываем платформу (tensorrt\_plan), названия и размерности входных и выходных тензоров. Описав параметры модели мы можем запустить данный Triton репозиторий непосредственно в самом Deepstream или запустить отдельный Triton сервер и подключить его к Deepstream через grpc.

Как описано выше, Nvidia Deepstream не сможет по умолчанию использовать такой формат выходных данных. Поэтому нужно создать объект прослойку, который будет считывать выходные тензоры и создавать объекты ограничивающих линий.

Для этого создадим класс NvInferServerCustomProcess, унаследовав его от класса IInferCustomProcessor. Этот класс отвечает за создание процесса, который может делать преобразование входных данных перед передачей в нейронную сеть и преобразование выходных данных на выходе с нейронной сети. За предобработку отвечает функция extraInputProcess, в нашем случае мы не будем делать никаких дополнительных преобразований при предобработке. 

Пост обработка описывается в функции inferenceDone. В начале нужно получить доступ к буферам выходных тензоров и проверить их размерности, как показано на рисунке ~\ref{src:src2_4}.
\begin{figure}
\begin{lstlisting}
auto num_detections = tensors["num"];
auto classes = tensors["classes"];
auto boxes = tensors["boxes"];
auto scores = tensors["scores"];
INFER_ASSERT(topDesc.dims.numDims == 2 && topDesc.dims.d[1] == 1);
INFER_ASSERT(boxDesc.dims.numDims == 3 && boxDesc.dims.d[2] == 4);
int32_t* classesPtr = (int32_t*)classes->getBufPtr(0);
int32_t* topPtr = (int32_t*)num_detections->getBufPtr(0);
float* boxesPtr = (float*)boxes->getBufPtr(0);
float* scoresPtr = (float*)scores->getBufPtr(0);
\end{lstlisting}
\caption{Получения буферов}
\label{src:src2_4}
\end{figure}

Теперь нужно обойти все объекты во всех батчах, для каждого из них создать объект NvDsInferObjectDetectionInfo и заполнить поля: classId, left, top, height, width, detectionConfidence. На выходе из модели у нас формат ограничивающего прямоугольника XYXY, а не XYWH, поэтому также нужно перевести из одного формата в другой. Код для превода в данный формат показан на рисунке ~\ref{src:src2_5}.

\begin{figure}
\begin{lstlisting}
obj.classId = (uint32_t)classesPtr[(i*100 + j)];
float* bbox = &boxesPtr[4*(i*100 + j)];
obj.left =  bbox[0];
obj.top = bbox[1];
float right = bbox[2];
float bottom = bbox[3];
obj.height = bottom - obj.top;
obj.width = right - obj.left;
obj.detectionConfidence = (float)scoresPtr[(i*100 + j)];
\end{lstlisting}
\caption{Получение данных}
\label{src:src2_5}
\end{figure}

Осталось только создать сами объекты для хранения данных объектов в Deepstream — objMeta и прикрепить их к самим кадрам видео. Из особенностей этого процесса стоит отметить, что необходимо привести координаты и размеры к размерам кадра до обработки (по умолчанию Yolo работаем с изображениями размера 640x640). Код для нормализации кадра приведен на рисунке ~\ref{src:src2_6}.

\begin{figure}
\begin{lstlisting}
rect_params.left = (obj.left * surfParamsList[batchIdx]->width) / 640.0;
rect_params.width = (obj.width * surfParamsList[batchIdx]->width) / 640.0;
rect_params.top = (obj.top * surfParamsList[batchIdx]->height) / 640.0;
rect_params.height = (obj.height * surfParamsList[batchIdx]->height) / 640.0;
\end{lstlisting}
\caption{Нормализация размера}
\label{src:src2_6}
\end{figure}

% Для распознавания номеров, был взят за основу фреймворк NomeroffNet, по умолчанию этот фреймворк тратит по 115 миллисекунд на обработку одного кадра, что никак не подходит для работы в реальном времени. В виде единого фреймворка для python нельзя сильно его ускорить и интегрировать в конвейер. Поэтому необходимо разбить данный вреймворк на участки препроцессинга и постпроцессинга данных между моделями и сами модели. 

% Изначально в фреймворке все процессы были представлены в виде унаследованных друг от друга классов, в связи с этим разбить работу библиотеки отдельно на преобразование данных и отдельно на применения моделей заняло некоторое время. Также часть моделей скачивалось в формате Pytorch при первом использовании библиотеки, а одна модель была представлена как составная модель на Pytorch Lightning, где одна часть это модель Pytorch, а вторая это чекпоинт использующий эту модель написанный на Pytorch Lightning, веса эти двух частей скачивались отдельно. Из-за того, что модель на Pytorch Lightning включала в себя части написанные на Pytorch, не удалось предварительно ее конвертиковать в формат Pytorch JIT, как остальные модели, а пришлось сразу конвертировать в формат ONNX.

% На вход этой системы будет приходить вырезанное изображение автомобиля, а на выходе необходимо получить текст государственного регистрационного номера. Сам конвейер будет состоять из нейросети Yolov5 обученной на детекцию номеров, которую была сконвертирована методом описанным выше. Также были применены модели Craft и Craft Refiner для построения точной границы текста и модель для распознавания текста на выровенно и нормированном изображении. Между моделями были написаны скрипты предобрабработки и постобработки данных.

Для обеспечения эффективного распознавания номеров в режиме реального времени основой послужила платформа NomeroffNet \cite{Nomeroffnet}. Однако время обработки одного кадра в этом фреймворке по умолчанию, которое составляет 115 миллисекунд, делает его непригодным для приложений реального времени. Следовательно, возникает необходимость разделить фреймворк на отдельные части: предварительная обработка данных, постобработка данных и сами модели, чтобы достичь желаемого ускорения и плавной интеграции в конвейере.

Изначально все процессы в рамках фреймворка были инкапсулированы как классы, унаследованные друг от друга. Следовательно, потребовались некоторые усилия, чтобы разделить функциональные возможности библиотеки на отдельные компоненты, особенно для преобразования данных и применения моделей. Кроме того, во время первоначального использования библиотеки было получено несколько моделей в формате PyTorch. Среди этих моделей одна была представлена как составная модель в PyTorch Lightning, где одна часть состояла из модели PyTorch, а другая часть была чекпоинтом, использующей эту модель в PyTorch Lightning. Веса для этих двух компонентов были загружены отдельно. Поскольку модель PyTorch Lightning включала в себя части, написанные на PyTorch, было невозможно напрямую преобразовать ее в формат PyTorch JIT, как в других моделях. Вместо этого потребовалось сразу преобразование в формат ONNX.

Эта система ожидает получения обрезанного изображения автомобиля в качестве входных данных и стремится извлечь текст номерного знака в качестве выходных данных. Сам конвейер включает в себя нейронную сеть Yolov5, которая была обучена распознавать числа. Процесс преобразования, описанный выше, был применен к модели Yolov5. Кроме того, для точного выделения областей текста на выровненном и нормализованном изображении были использованы модели Craft и Craft Refiner. Для эффективного соединения этих моделей были реализованы сценарии предварительной и постобработки.

Для работы библиотек, которые используются в Python-скриптах, необходимо сформировать и экспортировать окружение. Для этого требуется создать соответствующее окружение в Conda, а затем экспортировать его согласно указаниям, приведенным на изображении ~\ref{src:src2_7}.

\begin{figure}
\begin{lstlisting}
conda create --name pyenv python=3.8
conda activate pyenv
python -m pip install -r python_models/requirements.txt
conda-pack  
\end{lstlisting}
\caption{Упаковка окружения}
\label{src:src2_7}
\end{figure}

Файл “pyenv.tar.gz ” это упакованная среда, которая позже будет использоваться в Triton. Затем нам нужно скопировать файл во все модели python. После того, как среда будет готова, нам нужно указать Triton использовать эту среду для моделей, установив параметр <<EXECUTION\_ENV\_PATH>> в конфигурационных файлах, он должен быть уже указан в предоставленных конфигурациях. 

Чтобы подключить python код в качестве моделей необходимо описать вызов функций в методе execute класса TritonPythonModel, там будут описаны имена входных и выходных тензеров для python кода. Таким образом Triton будет рассматривать этот код как нейросетевую модель. На рисунке \ref{src:src3_3} приведен пример описание такого класса .

\begin{figure}
\begin{lstlisting}
import triton_python_backend_utils as pb_utils
class TritonPythonModel:
    def execute(self, requests):
        responses = []
        for request in requests:
            input_tensor_names = \
                ['bboxes', 'classes', 'masks', 'scores', 'shape']
            inputs = {
                tensor_name: pb_utils.get_input_tensor_by_name(request, tensor_name).as_numpy()
                      for tensor_name in input_tensor_names
            }
            predictions = self.postprocess(predictions)
            out_tensors = []
            for name in ['bboxes', 'classes', 'scores']:
                tensor = pb_utils.Tensor('post_' + name, predictions[name])
                out_tensors.append(tensor)
            response = pb_utils.InferenceResponse(output_tensors=out_tensors)
            responses.append(response)
        return responses
\end{lstlisting}
\caption{Пример описание python кода как нейросетевой модели для Nvidia Triton}
\label{src:src3_3}
\end{figure}


% Теперь нужно описать конвейер моделей в Nvidia Triton. Для этого необходимо подключить и описать конфигуроционные файлы в Nvidia Triton для всех моделей и скриптов по отдельности. Там нужно указать размерности и названия входных и выходных тензоров. Далее в той же папке репозитория где лежат и остальные модели нужно создать папку, где будет только описан конфигурационный файл для объединения других моделей. Для этого необходимо придумать названия входных,входных и промежуточных тензоров. Далее необходимо давать эти названия одной модели как входные тензоры, а другой как выходные, таким образом мы будем передавать данные от одной модели к другой.

% Далее необходимо подключить полученную единую модель к Nvidia Deepstream. Распознавание номеров будет описано как классификатор. Для этого также как и с Yolo необходимо создать динамическую библиотеку для записи номера как наименования класса. Nvidia Deepstream не может просто считать строку из тензора, поэтому необходимо записывать их в выходной тензор как числа в формате ASCII, а затем в Nvidia Deepstream перевести их обратно в формат строки. Ниже приведен пример конфигурационного файла данного конвейера.


Теперь необходимо описать конвейер моделей в Nvidia Triton. Для этого требуется подключить и описать конфигурационные файлы в Nvidia Triton для каждой модели и скрипта отдельно. В конфигурационных файлах следует указать размерности и названия входных и выходных тензоров. Далее, в той же папке репозитория, где находятся остальные модели, необходимо создать отдельную папку, в которой будет содержаться конфигурационный файл, описывающий объединение других моделей. При создании этого файла следует придумать названия входных, выходных и промежуточных тензоров. Затем следует указать эти названия входными тензорами для одной модели и выходными тензорами для другой, таким образом осуществляется передача данных от одной модели к другой.

После этого необходимо подключить полученную объединенную модель к Nvidia Deepstream. Распознавание номеров будет представляться в виде классификатора. Для этого, аналогично Yolo, требуется создать динамическую библиотеку, которая будет записывать номера в качестве названий классов. Поскольку Nvidia Deepstream не может просто считать строку из тензора, необходимо записывать номера в выходной тензор в виде ASCII-кодов чисел, а затем в Nvidia Deepstream переводить их обратно в строковый формат. На рисунке \ref{src:src2_8} приведен пример конфигурационного файла для данного конвейера.

\begin{figure}
\begin{lstlisting}
name: "pipeline"
platform: "ensemble"
max_batch_size: 0
input [{name: "IMAGE_BYTES"
    data_type: TYPE_FP32
    dims: [-1 ,640, 640,3]}]
output [{ name: "ZONE"
    data_type: TYPE_FP32
    dims: [1, 4]},
  {name: "TEXT"
    data_type: TYPE_INT32
    dims: [1, 9]}]
ensemble_scheduling {
  step [{model_name: "craft"
      model_version: -1
      input_map: {
        key: "image_bytes"
        value: "IMAGE_BYTES"}
      output_map: {
        key: "preprocessed_image"
        value: "PREPROCESSED_IMAGE"}
    {model_name: "car_postprocess"
      model_version: 1
      input_map: {
        key: "y"
        value: "PREPROCESSED_IMAGE""}
      output_map: {
        key: "text"
        value: "TEXT"}}]}

\end{lstlisting}
\caption{Пример конвейера в Nvidia Triton}
\label{src:src2_8}
\end{figure}

Таким образом общий конвейер представлен на рисунке \ref{fig:fig03}.

\begin{figure}
  \includegraphics[width=16cm]{inc/arch.png}
  \caption{Общий конвейер}
  \label{fig:fig03}
\end{figure}


